\documentclass[11pt,letterpaper]{article}

% ============================================================
% PAQUETES
% ============================================================
\usepackage[utf8]{inputenc}
\usepackage[T1]{fontenc}
\usepackage[spanish,es-noindentfirst,es-tabla]{babel}
\decimalpoint
\usepackage[margin=2.5cm]{geometry}
\usepackage{graphicx}
\usepackage{booktabs}
\usepackage{array}
\usepackage{multirow}
\usepackage{siunitx}
\sisetup{output-decimal-marker={.}, group-digits=false, detect-weight=true, table-number-alignment=center}
\usepackage{amsmath,amssymb}
\usepackage{listings}
\usepackage{xcolor}
\usepackage{fancyhdr}
\usepackage{titlesec}
\usepackage{caption}
\usepackage[hidelinks]{hyperref}

\graphicspath{{figuras/}}
\renewcommand{\topfraction}{0.92}
\renewcommand{\bottomfraction}{0.85}
\renewcommand{\textfraction}{0.06}
\renewcommand{\floatpagefraction}{0.80}
\setcounter{topnumber}{3}
\setcounter{totalnumber}{5}
\makeatletter
\setlength{\@fptop}{0pt}
\setlength{\@fpsep}{12pt}
\makeatother
\captionsetup{font=small,labelfont=bf}

% ============================================================
% ESTILO DE CÓDIGO
% ============================================================
\lstdefinestyle{codigo}{
    basicstyle=\ttfamily\footnotesize,
    keywordstyle=\color{blue!60!black}\bfseries,
    commentstyle=\color{gray},
    stringstyle=\color{orange!70!black},
    numbers=left,
    numberstyle=\tiny\color{gray},
    frame=single,
    breaklines=true,
    showstringspaces=false,
    tabsize=2,
    backgroundcolor=\color{gray!5},
    extendedchars=true,
    columns=fullflexible,
    keepspaces=true,
    upquote=true,
    literate={á}{{\'a}}1 {é}{{\'e}}1 {í}{{\'i}}1 {ó}{{\'o}}1 {ú}{{\'u}}1
             {Á}{{\'A}}1 {É}{{\'E}}1 {Í}{{\'I}}1 {Ó}{{\'O}}1 {Ú}{{\'U}}1
             {ñ}{{\~n}}1 {Ñ}{{\~N}}1 {¿}{{?`}}1 {¡}{{!`}}1
}
\lstset{style=codigo}

% ============================================================
% ENCABEZADO / PIE DE PÁGINA
% ============================================================
\pagestyle{fancy}
\fancyhf{}
\lhead{Taller 2 --- MLP y Backpropagation}
\rhead{Inteligencia Computacional Aplicada}
\cfoot{\thepage}
\renewcommand{\headrulewidth}{0.4pt}

\titleformat{\section}{\normalfont\Large\bfseries}{\thesection.}{0.5em}{}
\titleformat{\subsection}{\normalfont\large\bfseries}{\thesubsection}{0.5em}{}

% ============================================================
% DATOS DEL INFORME
% ============================================================
\newcommand{\nombreEstudiante}{Sergio Santiago Duarte Rojas}
\newcommand{\codigoEstudiante}{20261595002}
\newcommand{\nombreDocente}{Cesar Andrey Perdomo Charry}
\newcommand{\nombreAsignatura}{Inteligencia Computacional Aplicada}
\newcommand{\fechaInforme}{\today}

\begin{document}

% ============================================================
% PORTADA
% ============================================================
\begin{titlepage}
    \centering
    \vspace*{2cm}
    {\LARGE \bfseries Taller 2: Redes Neuronales Multicapa (MLP)\\ y Backpropagation\par}
    \vspace{0.5cm}
    {\large Informe de Resultados\par}
    \vspace{2cm}

    \begin{tabular}{ll}
        \textbf{Asignatura:} & \nombreAsignatura \\
        \textbf{Docente:}    & \nombreDocente \\
        \textbf{Estudiante:} & \nombreEstudiante \\
        \textbf{Código:}     & \codigoEstudiante \\
        \textbf{Fecha:}      & \fechaInforme \\
    \end{tabular}

    \vfill
    Universidad Distrital Francisco José de Caldas
\end{titlepage}

\tableofcontents
\newpage

% ============================================================
\section{Introducción}

El Taller 1 cerró con una limitación que ninguno de los dos modelos de una sola neurona pudo superar. El Perceptrón resolvió las compuertas AND y OR y el Adaline las aproximó con a lo sumo un patrón mal clasificado, pero frente a la XOR ambos quedaron en el 25 por ciento de acierto con la semilla usada allí, porque una neurona solo puede trazar un hiperplano y la XOR necesita dos. Este taller retoma exactamente ese punto. El objetivo es implementar un perceptrón multicapa, con al menos una capa oculta y entrenado por retropropagación del error, y comprobar sobre problemas de dificultad creciente qué gana la red al agregar esa capa.

El trabajo se desarrolló sobre el \textit{live script} base entregado por el docente, en el cual se completaron las cuatro funciones marcadas como pendientes sin alterar sus nombres ni sus argumentos. A partir de ahí se agregaron los experimentos de los numerales 3 a 6 del enunciado. Primero la XOR de dos y tres entradas, comparada contra los resultados del Taller 1. Después el problema multiclase de Iris con varias arquitecturas y tasas de aprendizaje. Luego los conjuntos Wine y Breast Cancer Wisconsin en las cuatro proporciones de partición, con el Perceptrón y el Adaline como referencia, y al final el análisis del sobreajuste y la parada temprana que lo mitiga. El informe cierra con la discusión de las cuatro preguntas que plantea la plantilla y con las conclusiones que se derivan de los datos.

% ============================================================
\section{Marco teórico (resumen)}

\subsection{Arquitectura del perceptrón multicapa}

Un perceptrón multicapa organiza las neuronas en capas sucesivas. La capa de entrada solo distribuye el patrón, cada capa oculta calcula una suma ponderada de las salidas de la capa anterior y le aplica una función de activación, y la capa de salida entrega la respuesta de la red. Para la neurona $j$ de la capa $k$ el cálculo es el de la ecuación \eqref{eq:forward}, donde $a^{(0)}$ es la entrada $x$.
\begin{equation}
    net_j^{(k)} = \sum_i w_{ji}^{(k)} a_i^{(k-1)} + b_j^{(k)}, \qquad a_j^{(k)} = f\left(net_j^{(k)}\right)
    \label{eq:forward}
\end{equation}
Este recorrido desde la entrada hasta la salida se llama propagación hacia adelante. La diferencia con la neurona única del Taller 1 está en que la función $f$ de las capas ocultas debe ser derivable y no lineal, con la salvedad de la ReLU, que no es derivable en cero y a la que ahí se le asigna derivada cero. El taller ofrece tres opciones, la sigmoidal $1/(1+e^{-net})$ que produce valores entre cero y uno, la tangente hiperbólica que produce valores entre menos uno y uno, y la ReLU, que vale $net$ cuando es positivo y cero en caso contrario.

\subsection{El algoritmo de retropropagación}

La red se entrena minimizando el error cuadrático $E = \tfrac{1}{2}\sum (d(x) - Y(x))^2$ por descenso de gradiente, igual que el Adaline. La novedad es que ahora hay pesos en capas que no ven directamente el error, y para corregirlos hace falta repartir ese error hacia atrás. El reparto se hace con un término $\delta$ por neurona. En la capa de salida el término se calcula con el error observado, ecuación \eqref{eq:delta_out}, y en cada capa oculta se obtiene a partir de los $\delta$ de la capa siguiente pesados por las conexiones que las unen, ecuación \eqref{eq:delta_hid}.
\begin{align}
    \delta^{(L)} &= \left(d(x) - Y(x)\right) \odot f'\left(net^{(L)}\right) \label{eq:delta_out} \\
    \delta^{(k)} &= \left(W^{(k+1)T}\delta^{(k+1)}\right) \odot f'\left(net^{(k)}\right) \label{eq:delta_hid}
\end{align}
Con los $\delta$ calculados, cada peso se corrige en proporción al $\delta$ de la neurona que alimenta y a la activación que recibe, ecuación \eqref{eq:update}. El coeficiente de momento $\beta$, que la guía de laboratorio del curso propone variar entre 0 y 1, agrega una fracción del incremento anterior, lo cual acelera el avance cuando los incrementos sucesivos apuntan en la misma dirección.
\begin{equation}
    \Delta W^{(k)} = \eta\, \delta^{(k)} \left(a^{(k-1)}\right)^T + \beta\, \Delta W^{(k)}_{\text{anterior}}, \qquad W^{(k)} \leftarrow W^{(k)} + \Delta W^{(k)}
    \label{eq:update}
\end{equation}
Cuando la red tiene una sola capa y activación lineal, las ecuaciones \eqref{eq:delta_out} y \eqref{eq:update} se reducen a la Regla Delta del Taller 1. La retropropagación es entonces la extensión natural de ese algoritmo a redes con capas ocultas.

\subsection{Por qué la XOR necesita una capa oculta}

Una neurona con función escalón divide el espacio de entrada en dos regiones separadas por un hiperplano. La XOR de dos entradas asigna la clase uno a los puntos $(0,1)$ y $(1,0)$ y la clase cero a $(0,0)$ y $(1,1)$, y esos cuatro puntos no admiten ninguna recta que deje a un lado los de una clase y al otro los de la otra. Por eso el Perceptrón y el Adaline del Taller 1 no pudieron resolverla. Rumelhart, Hinton y Williams (1986) mostraron que una capa oculta entrenada por retropropagación resuelve exactamente este tipo de problema. Con una capa oculta la red puede trazar dos rectas y combinarlas. Dos neuronas ocultas bastan en principio, una que se activa cuando alguna entrada vale uno y otra que se activa cuando ambas valen uno, y la neurona de salida aprende a responder uno solo cuando la primera está activa y la segunda no. La capa oculta transforma el problema en otro que sí es linealmente separable.

\subsection{Sobreajuste y parada temprana}

Una red con suficientes pesos y suficientes épocas puede aprender los datos de entrenamiento de memoria, incluidos su ruido y sus casos particulares. El síntoma es que el error de entrenamiento sigue bajando mientras el error sobre datos no vistos deja de bajar y empieza a subir. Ese fenómeno se llama sobreajuste (Haykin, 2009). La parada temprana lo controla vigilando el error sobre un conjunto de validación durante el entrenamiento y deteniéndolo cuando ese error lleva cierto número de épocas sin mejorar, conservando los pesos del mejor punto. El criterio de parada por épocas sin mejora que pide el numeral 2 del enunciado es exactamente ese mecanismo.

% ============================================================
\section{Metodología}

\subsection{Implementación}

Se completaron las cuatro funciones pendientes de la plantilla. \texttt{init\_mlp} construye las matrices de pesos de cada capa con valores aleatorios pequeños y los sesgos en cero, \texttt{forward\_mlp} aplica la ecuación \eqref{eq:forward} capa por capa y guarda las sumas ponderadas y las activaciones que necesita el paso siguiente, \texttt{backward\_mlp} calcula los $\delta$ de las ecuaciones \eqref{eq:delta_out} y \eqref{eq:delta_hid} desde la salida hacia la entrada y actualiza los pesos con la ecuación \eqref{eq:update}, y \texttt{activation\_deriv} entrega la derivada de cada función de activación. La activación de las capas ocultas se indica con un solo nombre para todas o con una lista con una por capa, de modo que una misma red puede combinar tangente hiperbólica y sigmoide. El código de las cuatro se reproduce en el Anexo \ref{anexo:codigo}. La plantilla usaba la estructura \texttt{net} en el ciclo de entrenamiento sin haberla creado, así que se agregó la llamada a \texttt{init\_mlp} en la sección correspondiente.

Para los experimentos se agregó una función \texttt{train\_mlp} que repite el ciclo de épocas de la plantilla y además calcula en cada época el error sobre un conjunto de validación, con lo cual implementa la parada por épocas sin mejora a través del campo \texttt{patience} de la configuración. El resto de las funciones agregadas son utilidades, la predicción de clase por el umbral de la configuración o por mayor salida, la carga de cada conjunto, la partición con semilla y el escalado al rango cero a uno, estas dos últimas tomadas del Taller 1. La partición en entrenamiento y prueba de Wine y Breast Cancer se entrega además como el script \texttt{Dataset\_train\_test\_T2}, versión análoga del script de clase del Taller 1 que lee los dos archivos de UCI y corta en las cuatro proporciones. El Perceptrón y el Adaline del Taller 1 se incluyeron como funciones locales sin cambios, y para los problemas de más de dos clases se entrena una neurona por clase y se decide por la mayor suma ponderada.

\subsection{Arquitecturas y parámetros evaluados}

La Tabla \ref{tab:config} recoge las configuraciones probadas. Todas usan activación sigmoidal en la salida, pesos iniciales de escala 0.5 y una tasa de aprendizaje de 0.3 salvo donde se indica. En las compuertas cada configuración se repite con cinco semillas y se reporta el promedio. Las arquitecturas escritas entre corchetes tienen dos capas ocultas. El error de parada por objetivo se dejó en el valor de la plantilla, 0.001. El Perceptrón y el Adaline conservan en todos los experimentos los parámetros del Taller 1, tasa 0.1 y 100 épocas el primero y tasa 0.05 y 200 épocas el segundo. El error que reportan las tablas y figuras como MSE, o error cuadrático medio, es el promedio sobre los patrones de la mitad de la suma de errores al cuadrado, el mismo $E$ de la sección 2.2.

\begin{table}[!htb]
    \centering
    \small
    \caption{Configuraciones de MLP evaluadas. La parada es por error objetivo o por número de épocas sin mejora en validación}
    \label{tab:config}
    \scriptsize
    \setlength{\tabcolsep}{4pt}
    \begin{tabular}{llcccc l}
        \toprule
        \textbf{Experimento} & \textbf{Ocultas} & \textbf{Activación oculta} & $\boldsymbol{\eta}$ & \textbf{Momento} & \textbf{Épocas} & \textbf{Parada} \\
        \midrule
        XOR & 2, 4, 8, [4 4] & sigmoid, tanh, relu & 0.3 & 0, 0.5, 0.9 & 5000 & error 0.001 \\
        Iris & 3, 6, 12, [8 4] & sigmoid & 0.05, 0.1, 0.3 & 0 & 1000 & error 0.001 \\
        Conjuntos reales & 8 & sigmoid & 0.3 & 0 & 500 & 30 sin mejora \\
        Sobreajuste & 2, 8, 32, 64 & sigmoid & 0.1 & 0 & 200 a 1500 & ninguna, 50 sin mejora \\
        \bottomrule
    \end{tabular}
\end{table}

\subsection{Datasets y particiones}

La XOR de $n$ entradas se genera como tabla de verdad completa con salida uno cuando un número impar de entradas vale uno, con 4 patrones para dos entradas y 8 para tres. Iris tiene 150 flores, cuatro medidas y tres especies de cincuenta ejemplares, y se codifica con una salida por especie. Wine, del repositorio UCI, tiene 178 vinos de tres bodegas descritos por trece medidas químicas. Breast Cancer Wisconsin Diagnostic, también de UCI, tiene 569 tumores con treinta medidas del núcleo celular y dos clases, 357 benignos y 212 malignos. El conjunto de billetes es el mismo del Taller 1, con 1372 patrones, cuatro medidas y dos clases. La Tabla \ref{tab:tamanos} resume los tamaños.

\begin{table}[!htb]
    \centering
    \small
    \caption{Conjuntos de datos empleados}
    \label{tab:tamanos}
    \begin{tabular}{l S[table-format=4.0] S[table-format=2.0] S[table-format=1.0] l}
        \toprule
        \textbf{Conjunto} & {\textbf{Patrones}} & {\textbf{Entradas}} & {\textbf{Clases}} & \textbf{Uso} \\
        \midrule
        XOR de 2 entradas & 4 & 2 & 2 & numeral 3 \\
        XOR de 3 entradas & 8 & 3 & 2 & numeral 3 \\
        Iris & 150 & 4 & 3 & numeral 4 \\
        Wine & 178 & 13 & 3 & numeral 5 \\
        Breast Cancer Wisconsin & 569 & 30 & 2 & numerales 5 y 6 \\
        Billetes & 1372 & 4 & 2 & numeral 5 \\
        \bottomrule
    \end{tabular}
\end{table}

Las particiones se generan con la misma función del Taller 1, que mezcla los índices con una semilla fija y corta en la proporción pedida, 60-40, 70-30, 80-20 y 90-10. Las entradas se escalan al rango cero a uno con los mínimos y máximos del conjunto de entrenamiento, y esas mismas cotas se aplican al de prueba. Cuando el entrenamiento necesita un conjunto de validación, sea para la curva de error o para la parada temprana, se aparta una quinta parte del conjunto de entrenamiento con una segunda semilla, de modo que la exactitud siempre se mide sobre una prueba que la red no vio durante el entrenamiento. El Perceptrón y el Adaline, que no usan validación, se entrenan con el conjunto de entrenamiento completo. El script \texttt{Dataset\_train\_test\_T2} usa la misma semilla que el experimento y produce exactamente los mismos cortes de entrenamiento y prueba. Wine proviene de Aeberhard y Forina (1991), Breast Cancer de Wolberg et al. (1993) y billetes de Lohweg (2013), todos del repositorio UCI, e Iris es el conjunto clásico de Fisher (1936).

% ============================================================
\section{Resultados}

\subsection{Problema XOR}

La Tabla \ref{tab:xor} compara los tres modelos sobre la XOR de dos y tres entradas. El Perceptrón y el Adaline no convergen en ninguna semilla y se quedan en el 30 por ciento con dos entradas y cerca del 50 por ciento con tres, es decir por debajo o al nivel de responder siempre la misma clase, porque los pesos siguen oscilando cuando se agotan las épocas. El MLP con cuatro u ocho neuronas ocultas resuelve la XOR de dos entradas en todas las semillas, aunque necesita entre 3443 y 4117 épocas para llevar el error por debajo de 0.001 con la tasa base y sin momento. Con dos neuronas ocultas, el mínimo teórico, las cinco semillas clasifican bien los cuatro patrones, pero una no alcanza el error objetivo en 5000 épocas. La XOR de tres entradas resulta más exigente y con la misma configuración solo la capa de ocho neuronas converge en la mayoría de las semillas. La última fila de cada bloque adelanta el efecto del momento, que con 0.9 lleva a las cinco semillas al error objetivo en ambos problemas.

\begin{table}[!htb]
    \centering
    \small
    \caption{XOR de dos y tres entradas. Promedio de cinco semillas. El Perceptrón y el Adaline usan los parámetros del Taller 1, el MLP la tasa 0.3 sin momento}
    \label{tab:xor}
    \begin{tabular}{llcc S[table-format=4.1] S[table-format=3.0] S[table-format=3.1]}
        \toprule
        \textbf{Entradas} & \textbf{Modelo} & \textbf{Ocultas} & \textbf{Momento} & {\textbf{Épocas}} & {\textbf{Convergió (\%)}} & {\textbf{Exactitud (\%)}} \\
        \midrule
        \multirow{6}{*}{2} & Perceptrón & ninguna & {--} & 100.0  & 0   & 30.0  \\
                           & Adaline    & ninguna & {--} & 200.0  & 0   & 30.0  \\
                           & MLP        & 2       & 0    & 4406.0 & 80  & 100.0 \\
                           & MLP        & 4       & 0    & 4116.6 & 100 & 100.0 \\
                           & MLP        & 8       & 0    & 3443.0 & 100 & 100.0 \\
                           & MLP        & 4       & 0.9  & 435.8  & 100 & 100.0 \\
        \midrule
        \multirow{6}{*}{3} & Perceptrón & ninguna & {--} & 100.0  & 0   & 47.5  \\
                           & Adaline    & ninguna & {--} & 200.0  & 0   & 50.0  \\
                           & MLP        & 2       & 0    & 5000.0 & 0   & 62.5  \\
                           & MLP        & 4       & 0    & 4187.8 & 40  & 75.0  \\
                           & MLP        & 8       & 0    & 3122.8 & 80  & 92.5  \\
                           & MLP        & 4       & 0.9  & 1231.4 & 100 & 100.0 \\
        \bottomrule
    \end{tabular}
\end{table}

El momento y la función de activación cambian ese panorama, como muestra la Tabla \ref{tab:xor_param}. Con cuatro neuronas ocultas y momento 0.9 la XOR de dos entradas converge en 436 épocas en lugar de 4117, casi diez veces más rápido, y la de tres entradas pasa de converger en dos de cinco semillas a converger en todas. La tangente hiperbólica también acelera, con 880 épocas en dos entradas, mientras que la ReLU resulta poco fiable en este problema, con tres de cinco semillas convergidas en dos entradas y ninguna en tres. La red de dos capas ocultas de cuatro neuronas, con tangente hiperbólica en la primera y sigmoide en la segunda, es la más rápida de todas con momento 0.9, 149 épocas en dos entradas y 260 en tres. Como esa red combina la segunda capa con la tangente hiperbólica y el momento, la tabla no permite separar cuánto aporta cada uno. La Figura \ref{fig:xor} muestra la curva de error del MLP con momento 0.9 frente a la del Adaline, que se estabiliza cerca de 0.139, por encima del piso de 0.125 que alcanza la mejor función lineal al responder 0.5 en los cuatro patrones.

\begin{table}[!htb]
    \centering
    \small
    \caption{Efecto del momento y de la activación oculta sobre la XOR, y una red de dos capas ocultas con activación distinta en cada una. Promedio de cinco semillas}
    \label{tab:xor_param}
    \footnotesize
    \begin{tabular}{lcll S[table-format=4.1] S[table-format=3.0] S[table-format=3.1]}
        \toprule
        \textbf{Entradas} & \textbf{Ocultas} & \textbf{Activación oculta} & \textbf{Momento} & {\textbf{Épocas}} & {\textbf{Convergió (\%)}} & {\textbf{Exactitud (\%)}} \\
        \midrule
        \multirow{6}{*}{2} & 4     & sigmoid        & 0   & 4116.6 & 100 & 100.0 \\
                           & 4     & sigmoid        & 0.5 & 2087.6 & 100 & 100.0 \\
                           & 4     & sigmoid        & 0.9 & 435.8  & 100 & 100.0 \\
                           & 4     & tanh           & 0   & 879.8  & 100 & 100.0 \\
                           & 4     & relu           & 0   & 2416.4 & 60  & 90.0  \\
                           & {[4 4]} & tanh, sigmoid & 0.9 & 149.0 & 100 & 100.0 \\
        \midrule
        \multirow{6}{*}{3} & 4     & sigmoid        & 0   & 4187.8 & 40  & 75.0  \\
                           & 4     & sigmoid        & 0.5 & 4281.2 & 20  & 95.0  \\
                           & 4     & sigmoid        & 0.9 & 1231.4 & 100 & 100.0 \\
                           & 4     & tanh           & 0   & 1516.0 & 80  & 90.0  \\
                           & 4     & relu           & 0   & 5000.0 & 0   & 87.5  \\
                           & {[4 4]} & tanh, sigmoid & 0.9 & 259.6 & 100 & 100.0 \\
        \bottomrule
    \end{tabular}
\end{table}

\begin{figure}[!htb]
    \centering
    \includegraphics[width=12cm]{fig02.png}
    \caption{Curva de error de entrenamiento sobre la XOR. El MLP con cuatro neuronas ocultas y momento 0.9 llega al error objetivo, el Adaline se estanca cerca del piso de una función lineal, 0.125}
    \label{fig:xor}
\end{figure}

\subsection{Clasificación multiclase en Iris}

La Tabla \ref{tab:iris} recoge las doce combinaciones de arquitectura y tasa de aprendizaje sobre la partición 70-30, con 84 flores de entrenamiento, 21 de validación y 45 de prueba. Ninguna configuración alcanza el error objetivo de 0.001 en mil épocas, porque con tres salidas y clases que se traslapan el error cuadrático de entrenamiento termina entre 0.005 y 0.017. La exactitud de prueba es de 97.8 por ciento en nueve de las doce configuraciones, que corresponde a una sola flor mal clasificada, una \textit{versicolor} tomada por \textit{virginica}, y de 95.6 por ciento con la tasa 0.05 en las tres redes de una capa, que a las mil épocas todavía no terminan de ajustar la frontera entre esas dos especies y fallan una flor más. La mejor configuración se elige por el error de validación al final del entrenamiento, sin mirar la prueba, y resulta ser la de doce neuronas con tasa 0.05. La Tabla \ref{tab:iris_conf} muestra su matriz de confusión, con dos \textit{versicolor} tomadas por \textit{virginica}. El menor error de validación no garantizó la mayor exactitud de prueba, porque sobre 45 flores la diferencia entre las configuraciones es de una sola flor. Las nueve configuraciones con 97.8 por ciento comparten una misma matriz, y la Tabla \ref{tab:iris_todas} del Anexo \ref{anexo:tablas} reúne las doce.

\begin{table}[!htb]
    \centering
    \small
    \caption{Iris sobre la partición 70-30. Mil épocas en todas las configuraciones, exactitud sobre las 45 flores de prueba}
    \label{tab:iris}
    \footnotesize
    \begin{tabular}{c S[table-format=1.2] S[table-format=1.3] S[table-format=1.4] S[table-format=1.4] S[table-format=3.1]}
        \toprule
        \textbf{Neuronas ocultas} & {$\boldsymbol{\eta}$} & {\textbf{Segundos}} & {\textbf{MSE entrenamiento}} & {\textbf{MSE validación}} & {\textbf{Exactitud (\%)}} \\
        \midrule
        3       & 0.05 & 2.431 & 0.0159 & 0.0302 & 95.6 \\
        3       & 0.1  & 2.452 & 0.0114 & 0.0375 & 97.8 \\
        3       & 0.3  & 2.617 & 0.0106 & 0.0620 & 97.8 \\
        6       & 0.05 & 3.254 & 0.0166 & 0.0305 & 95.6 \\
        6       & 0.1  & 2.921 & 0.0115 & 0.0348 & 97.8 \\
        6       & 0.3  & 2.724 & 0.0066 & 0.0444 & 97.8 \\
        12      & 0.05 & 3.171 & 0.0167 & 0.0287 & 95.6 \\
        12      & 0.1  & 3.003 & 0.0117 & 0.0347 & 97.8 \\
        12      & 0.3  & 2.968 & 0.0095 & 0.0591 & 97.8 \\
        {[8 4]} & 0.05 & 3.925 & 0.0109 & 0.0343 & 97.8 \\
        {[8 4]} & 0.1  & 3.990 & 0.0068 & 0.0436 & 97.8 \\
        {[8 4]} & 0.3  & 5.231 & 0.0053 & 0.0444 & 97.8 \\
        \bottomrule
    \end{tabular}
\end{table}

\begin{table}[!htb]
    \centering
    \small
    \caption{Matriz de confusión sobre Iris de la configuración con menor error de validación, doce neuronas ocultas y tasa 0.05, 45 flores de prueba}
    \label{tab:iris_conf}
    \begin{tabular}{l S[table-format=2.0] S[table-format=2.0] S[table-format=2.0]}
        \toprule
        \textbf{Clase real} & {\textbf{Predicha setosa}} & {\textbf{Predicha versicolor}} & {\textbf{Predicha virginica}} \\
        \midrule
        setosa     & 14 & 0  & 0  \\
        versicolor & 0  & 11 & 2  \\
        virginica  & 0  & 0  & 18 \\
        \bottomrule
    \end{tabular}
\end{table}

La Figura \ref{fig:iris} muestra las curvas de error de entrenamiento y de validación de las doce configuraciones. La tasa gobierna la forma de la curva mucho más que el número de neuronas. Con 0.05 el error baja despacio y todavía desciende al llegar a mil épocas, con 0.1 se estabiliza cerca de la época 300, y con 0.3 llega antes pero oscila de forma visible durante todo el entrenamiento, porque cada patrón mueve los pesos demasiado. Con la tasa 0.3 el error de validación además se despega del de entrenamiento y sube despacio desde la época 200, señal de un sobreajuste leve sobre 84 patrones que no llega a costar flores de prueba. La red de dos capas ocultas alcanza el error de entrenamiento más bajo en las tres tasas y con 0.05 es la única que llega a 97.8 por ciento, pero con 0.1 y 0.3 no mejora la exactitud, porque las 45 flores ya se clasifican con una sola excepción.

\begin{figure}[!htb]
    \centering
    \includegraphics[width=15cm]{fig03.png}
    \caption{Curvas de error de entrenamiento y de validación sobre Iris para las cuatro arquitecturas y las tres tasas de aprendizaje}
    \label{fig:iris}
\end{figure}

\subsection{Wine y Breast Cancer Wisconsin en distintas particiones}

La Tabla \ref{tab:particiones} compara el MLP contra el Perceptrón y el Adaline sobre los tres conjuntos reales y las cuatro proporciones. Se agregó el conjunto de billetes porque es el que permite la comparación más directa con el Taller 1. Los valores del Perceptrón y del Adaline sobre él quedan cerca de los reportados allí y no coinciden exactamente, porque aquí el escalado usa las cotas de la parte de entrenamiento que no se reserva para validación y el orden de los patrones cambia. La Figura \ref{fig:particiones} muestra la exactitud de prueba de los tres modelos.

\begin{table}[!htb]
    \centering
    \small
    \caption{Exactitud de prueba según proporción de partición, con las épocas y el tiempo de entrenamiento en segundos. El MLP usa ocho neuronas ocultas y para por treinta épocas sin mejora en validación o por error objetivo, y sus épocas son las del mejor punto de validación o las de llegada al error objetivo. Las 200 épocas del Adaline y las 100 del Perceptrón en Breast Cancer y billetes son el tope, sin convergencia. Los tiempos cambian ligeramente en cada ejecución}
    \label{tab:particiones}
    \footnotesize
    \begin{tabular}{ll S[table-format=3.1] S[table-format=3.0] S[table-format=1.3] S[table-format=3.1] S[table-format=3.0] S[table-format=1.3] S[table-format=3.1] S[table-format=3.0] S[table-format=1.3]}
        \toprule
         & & \multicolumn{3}{c}{\textbf{MLP}} & \multicolumn{3}{c}{\textbf{Perceptrón}} & \multicolumn{3}{c}{\textbf{Adaline}} \\
        \cmidrule(lr){3-5} \cmidrule(lr){6-8} \cmidrule(lr){9-11}
        \textbf{Conjunto} & \textbf{Partición} & {\textbf{\%}} & {\textbf{Épocas}} & {\textbf{s}} & {\textbf{\%}} & {\textbf{Épocas}} & {\textbf{s}} & {\textbf{\%}} & {\textbf{Épocas}} & {\textbf{s}} \\
        \midrule
        \multirow{4}{*}{Wine}           & 60-40 & 98.6 & 119 & 0.605 & 95.8 & 52 & 0.006 & 97.2 & 200 & 0.025 \\
                                        & 70-30 & 98.1 & 150 & 0.594 & 92.5 & 32 & 0.004 & 96.2 & 200 & 0.018 \\
                                        & 80-20 & 100.0 & 118 & 0.555 & 94.4 & 25 & 0.007 & 97.2 & 200 & 0.022 \\
                                        & 90-10 & 94.4 & 121 & 0.612 & 100.0 & 35 & 0.003 & 94.4 & 200 & 0.023 \\
        \midrule
        \multirow{4}{*}{Breast Cancer}  & 60-40 & 97.4 & 52 & 0.743 & 96.1 & 100 & 0.019 & 95.2 & 200 & 0.018 \\
                                        & 70-30 & 97.7 & 64 & 1.091 & 95.3 & 100 & 0.011 & 97.1 & 200 & 0.028 \\
                                        & 80-20 & 95.6 & 72 & 1.225 & 93.0 & 100 & 0.013 & 96.5 & 200 & 0.023 \\
                                        & 90-10 & 98.2 & 88 & 1.523 & 98.2 & 100 & 0.013 & 98.2 & 200 & 0.024 \\
        \midrule
        \multirow{4}{*}{Billetes}       & 60-40 & 99.5 & 132 & 2.548 & 96.9 & 100 & 0.021 & 96.5 & 200 & 0.031 \\
                                        & 70-30 & 100.0 & 146 & 3.437 & 99.0 & 100 & 0.019 & 95.9 & 200 & 0.038 \\
                                        & 80-20 & 99.3 & 125 & 3.256 & 96.7 & 100 & 0.022 & 96.4 & 200 & 0.047 \\
                                        & 90-10 & 98.5 & 105 & 3.173 & 97.8 & 100 & 0.024 & 95.6 & 200 & 0.050 \\
        \bottomrule
    \end{tabular}
\end{table}

\begin{figure}[!htb]
    \centering
    \includegraphics[width=15cm]{fig04.png}
    \caption{Exactitud de prueba de los tres modelos según la proporción de partición en los tres conjuntos reales}
    \label{fig:particiones}
\end{figure}

Los tres conjuntos cuentan historias distintas. En Wine las tres bodegas resultan separables por hiperplanos y el Perceptrón converge en todas las particiones entre 25 y 52 épocas, y aun así el MLP lo supera en tres de las cuatro, con 98.1 a 100 por ciento frente a 92.5 a 95.8. Una explicación es que el Perceptrón se detiene en el primer hiperplano que separa el entrenamiento, sin importar lo cerca que quede de los patrones, mientras que el MLP sigue reduciendo el error hasta que la validación deja de mejorar. La excepción es la partición 90-10, donde el Perceptrón acierta los 18 vinos de prueba y el MLP falla uno. En Breast Cancer el Perceptrón nunca converge en cien épocas, señal de que las dos clases se traslapan, y el MLP lo supera en 60-40, 70-30 y 80-20, con 97.4, 97.7 y 95.6 por ciento frente a 96.1, 95.3 y 93.0, y lo iguala en 90-10 con 98.2. Frente al Adaline el MLP solo pierde en esa partición 80-20, 95.6 contra 96.5 por ciento, un tumor de 114. En billetes la ventaja del MLP es clara y consistente, entre 98.5 y 100 por ciento frente a un máximo de 99.0 del Perceptrón y 96.5 del Adaline. El Adaline queda entre los otros dos modelos en tres de las cuatro particiones de Wine, pero en billetes es el último en las cuatro, porque el hiperplano que minimiza el error cuadrático no es el que mejor separa las clases.

El costo de esa ventaja es el tiempo. El MLP tarda entre 0.555 y 3.437 segundos por entrenamiento contra menos de 0.1 segundos del Perceptrón y el Adaline, entre 25 y 194 veces más, porque cada patrón atraviesa dos capas hacia adelante y dos hacia atrás y además la red evalúa su error de validación en cada época. En cuanto a las épocas, en Wine y Breast Cancer el MLP alcanza su mejor punto de validación entre las épocas 52 y 150, y la parada llega treinta épocas después. En billetes, en cambio, el MLP no llega a usar la parada temprana, porque su error de entrenamiento baja del objetivo de 0.001 entre las épocas 105 y 146. El Adaline agota sus 200 épocas en los tres conjuntos y el Perceptrón sus 100 en Breast Cancer y billetes, así que en esos casos sus épocas y su tiempo son los del tope y no los de una convergencia. Sobre la proporción de partición no aparece una tendencia común, y las particiones 90-10 evalúan sobre 18, 57 y 137 patrones, así que cada error individual pesa entre 0.7 y 5.6 puntos y sus cifras son las menos estables.

\subsection{Análisis de sobreajuste (overfitting)}

La Figura \ref{fig:sobreajuste} muestra el error de entrenamiento y de validación sobre Breast Cancer durante 1500 épocas para cuatro anchos de capa oculta, sin parada temprana, y la Tabla \ref{tab:sobreajuste} resume las cifras. Las cuatro curvas tienen la misma forma. El error de validación baja hasta un mínimo entre 0.0068 y 0.0079 en las épocas 85 a 165 y a partir de ahí sube de manera sostenida hasta terminar entre 0.0187 y 0.0220, casi el triple, mientras el error de entrenamiento sigue bajando hasta 0.0036. Esa separación creciente entre las dos curvas es la evidencia del sobreajuste. La red sigue ajustándose a los 318 patrones de entrenamiento y pierde capacidad sobre los 80 de validación y los 171 de prueba.

\begin{figure}[!htb]
    \centering
    \includegraphics[width=15cm]{fig05.png}
    \caption{Error de entrenamiento y de validación sobre Breast Cancer para cuatro anchos de capa oculta. La línea discontinua marca la época de mínimo error de validación}
    \label{fig:sobreajuste}
\end{figure}

\begin{table}[!htb]
    \centering
    \small
    \caption{Sobreajuste en Breast Cancer, partición 70-30. Sin parada temprana el entrenamiento dura 1500 épocas, con parada temprana se detiene tras cincuenta épocas sin mejora en validación}
    \label{tab:sobreajuste}
    \scriptsize
    \begin{tabular}{c S[table-format=1.4] S[table-format=1.4] S[table-format=3.0] S[table-format=1.4] S[table-format=2.1] S[table-format=3.0] S[table-format=2.1]}
        \toprule
         & \multicolumn{5}{c}{\textbf{Sin parada temprana, 1500 épocas}} & \multicolumn{2}{c}{\textbf{Con parada temprana}} \\
        \cmidrule(lr){2-6} \cmidrule(lr){7-8}
        \textbf{Ocultas} & {\textbf{MSE ent.}} & {\textbf{MSE val.}} & {\textbf{Época mín.}} & {\textbf{MSE val. mín.}} & {\textbf{Exact. (\%)}} & {\textbf{Épocas}} & {\textbf{Exact. (\%)}} \\
        \midrule
        2  & 0.0038 & 0.0187 & 165 & 0.0068 & 94.7 & 215 & 97.7 \\
        8  & 0.0036 & 0.0189 & 128 & 0.0079 & 94.7 & 178 & 97.7 \\
        32 & 0.0037 & 0.0220 & 130 & 0.0079 & 94.2 & 180 & 97.7 \\
        64 & 0.0036 & 0.0206 & 85  & 0.0068 & 94.7 & 135 & 97.7 \\
        \bottomrule
    \end{tabular}
\end{table}

La Tabla \ref{tab:epocas} aísla el efecto del número de épocas con ocho neuronas ocultas. Cada entrenamiento más largo baja el error de entrenamiento y sube el de validación, y la exactitud de prueba cae de 95.9 por ciento a las 200 épocas a 94.7 a las 1000 y 1500, es decir dos tumores más mal clasificados por seguir entrenando.

\begin{table}[!htb]
    \centering
    \small
    \caption{Efecto del número de épocas con ocho neuronas ocultas sobre Breast Cancer, sin parada temprana}
    \label{tab:epocas}
    \begin{tabular}{S[table-format=4.0] S[table-format=1.4] S[table-format=1.4] S[table-format=2.1]}
        \toprule
        {\textbf{Épocas}} & {\textbf{MSE entrenamiento}} & {\textbf{MSE validación}} & {\textbf{Exactitud de prueba (\%)}} \\
        \midrule
        200  & 0.0071 & 0.0115 & 95.9 \\
        500  & 0.0051 & 0.0126 & 95.3 \\
        1000 & 0.0040 & 0.0146 & 94.7 \\
        1500 & 0.0036 & 0.0189 & 94.7 \\
        \bottomrule
    \end{tabular}
\end{table}

El criterio de parada temprana implementado vigila el error de validación en cada época, guarda los pesos del mejor punto y detiene el entrenamiento cuando pasan cincuenta épocas sin mejorar ese mínimo, devolviendo los pesos guardados. Con ese criterio el entrenamiento termina entre las épocas 135 y 215, cincuenta después del mínimo como corresponde, y la exactitud de prueba sube a 97.7 por ciento en los cuatro anchos, frente a valores entre 94.2 y 94.7 sin parada. Cada una de las cuatro redes detenidas falla cuatro de los 171 tumores de prueba. El número de neuronas ocultas, en cambio, apenas cambia el resultado. Con dos neuronas y con sesenta y cuatro el mínimo de validación vale lo mismo, y la red más ancha solo lo alcanza antes, de modo que en este conjunto el sobreajuste lo produce la duración del entrenamiento y no el tamaño de la red.

% ============================================================
\section{Análisis y discusión}

\subsection{Por qué el MLP resuelve la XOR y los modelos del Taller 1 no}

Los resultados de la Tabla \ref{tab:xor} reproducen el límite que cerró el Taller 1 y muestran cómo se supera. El Perceptrón y el Adaline quedan en 30 por ciento con dos entradas y cerca de 50 con tres, porque una sola neurona traza un solo hiperplano y ningún hiperplano separa los patrones de la XOR. El MLP con capa oculta transforma la entrada antes de decidir. Cada neurona oculta traza su propio hiperplano y la neurona de salida combina las regiones resultantes, así que el problema que llega a la salida ya es separable. La fila de dos neuronas ocultas, el mínimo teórico, confirma que dos rectas bastan, pues clasifica bien los cuatro patrones en las cinco semillas, aunque una de ellas no baja el error de 0.001 en 5000 épocas. Con cuatro u ocho neuronas sobran caminos y en dos entradas todas las semillas llegan al error objetivo.

La XOR de tres entradas explica por qué el problema crece rápido. Con tres entradas hay ocho vértices de un cubo alternando clase, y separarlos exige más regiones, así que con la tasa base solo la capa de ocho neuronas converge en la mayoría de las semillas. El momento resultó ser la herramienta decisiva en ese caso. Con $\beta$ igual a 0.9 la corrección acumula impulso en la dirección que se repite época tras época y atraviesa las mesetas donde el gradiente es pequeño, que en la Figura \ref{fig:xor} se ven como los tramos planos de la curva de tres entradas. El resultado es convergencia en todas las semillas y casi diez veces menos épocas en dos entradas.

\subsection{Efecto del número de neuronas ocultas y de la tasa de aprendizaje}

La tasa de aprendizaje mostró el mismo compromiso en todos los experimentos. Un valor pequeño produce curvas suaves pero lentas, como la tasa 0.05 en Iris, que a las mil épocas todavía no se ha estabilizado. Un valor grande llega antes al vecindario del mínimo pero oscila alrededor de él, como la tasa 0.3 en la Figura \ref{fig:iris}, y en Iris esa oscilación despega el error de validación del de entrenamiento. El valor intermedio de 0.1 dio la mejor combinación de velocidad y estabilidad, y el de 0.05 dejó la frontera entre \textit{versicolor} y \textit{virginica} sin terminar de ajustar a las mil épocas en las redes de una capa. En la XOR solo se probó la tasa 0.3, y con ella fue el momento lo que permitió avanzar con paso grande sin perder estabilidad.

El número de neuronas ocultas tuvo un efecto menor de lo que se podría esperar. En Iris las cuatro arquitecturas dan la misma exactitud con las tasas 0.1 y 0.3, y en Breast Cancer las cuatro anchuras probadas caen en el mismo error de validación mínimo, solo que la más ancha llega antes. Solo en la XOR, donde la geometría exige un número mínimo de hiperplanos, la anchura marca la diferencia entre converger o no. En los conjuntos reales las clases se separan con pocas neuronas y las adicionales no aportan capacidad útil, solo más pesos que ajustar y más tiempo por época, como muestra la columna de segundos de la Tabla \ref{tab:iris}.

\subsection{Proporción de partición y generalización}

Ninguna proporción resultó mejor de manera consistente. En Wine el Perceptrón alcanza su máximo en 90-10 y el MLP en 80-20, en Breast Cancer el MLP se mueve entre 95.6 y 98.2 por ciento sin orden con la proporción, y en billetes las cuatro particiones dan entre 98.5 y 100 por ciento. El patrón que sí se repite es la fragilidad de las particiones con conjuntos de prueba pequeños. Con 18 vinos o 57 tumores de prueba, un solo error cambia la exactitud en varios puntos, y por eso el 100 por ciento del Perceptrón en Wine 90-10 o del MLP en Wine 80-20, con 36 vinos de prueba, dicen menos de lo que parecen. Con una sola semilla por partición no fue posible identificar una proporción mejor que las demás.

\subsection{Evidencia de sobreajuste y su mitigación}

La evidencia está en la Figura \ref{fig:sobreajuste} y en la Tabla \ref{tab:epocas}. En los cuatro anchos el error de validación toca un mínimo entre las épocas 85 y 165 y después sube de forma sostenida, mientras el de entrenamiento baja sin parar. A las 1500 épocas la red clasifica peor los datos nuevos que a las 200, con una diferencia de 1.2 puntos de exactitud y de 2.9 a 3.5 puntos frente a la red detenida a tiempo, aunque su error sobre los datos que ya vio sea cinco veces menor que el de validación. Ese divorcio entre los dos errores es la definición operativa del sobreajuste.

La parada temprana lo mitigó en los cuatro casos y los llevó al mismo 97.7 por ciento. El criterio es sencillo y no exige ninguna decisión previa sobre el número de épocas, porque es el propio conjunto de validación el que la toma. Lo que sí exige es apartar datos para validar, que aquí fueron 80 de los 398 patrones de entrenamiento, con lo cual la red aprende con menos ejemplos pero la exactitud sobre la prueba mide de verdad lo que la red no vio. Un hallazgo que conviene subrayar es que aumentar las neuronas ocultas de dos a sesenta y cuatro no agravó el sobreajuste, pues el mínimo de validación y la exactitud final quedan en el mismo rango, aunque lo adelantó de la época 165 a la 85. La causa en este conjunto es el exceso de épocas, y el remedio correcto es detener el entrenamiento a tiempo y no reducir la red.

% ============================================================
\section{Conclusiones}

El perceptrón multicapa resolvió la XOR de dos entradas con exactitud del cien por ciento en todas las semillas con cuatro u ocho neuronas ocultas, y la de tres entradas con la misma exactitud al agregar momento 0.9, mientras que el Perceptrón y el Adaline del Taller 1 quedaron en 30 por ciento con dos entradas y entre 47.5 y 50 con tres. La capa oculta es la diferencia, y dos neuronas ocultas bastan para clasificar los cuatro patrones, aunque con ellas una de cinco semillas no llega al error objetivo.

El momento fue el parámetro de mayor efecto sobre la velocidad. Con $\beta$ igual a 0.9 la XOR de dos entradas pasó de 4117 a 436 épocas y la de tres entradas pasó de converger en dos semillas a converger en las cinco. La tangente hiperbólica también aceleró la convergencia y la ReLU resultó poco fiable en un problema tan pequeño.

Sobre Iris las doce configuraciones alcanzaron 95.6 o 97.8 por ciento de exactitud, y la tasa de aprendizaje determinó la forma de la curva de error más que el número de neuronas o de capas. La tasa 0.1 dio el mejor equilibrio entre velocidad y estabilidad.

En los conjuntos reales el MLP superó al Perceptrón en diez de las doce particiones y al Adaline en nueve, y solo perdió dos veces, en Wine 90-10 frente al Perceptrón y en Breast Cancer 80-20 frente al Adaline, cada vez por un solo patrón de prueba. La ventaja más clara estuvo en billetes, entre 98.5 y 100 por ciento. El precio fue un tiempo de entrenamiento entre 25 y 194 veces mayor.

El sobreajuste apareció en Breast Cancer como una subida sostenida del error de validación desde las épocas 85 a 165 mientras el de entrenamiento seguía bajando, costó 1.2 puntos de exactitud entre 200 y 1500 épocas, y su magnitud no dependió del ancho de la capa oculta entre dos y sesenta y cuatro neuronas, aunque la red más ancha llegó antes al mínimo. La parada temprana por cincuenta épocas sin mejora lo corrigió y elevó la exactitud de prueba de entre 94.2 y 94.7 a 97.7 por ciento en los cuatro casos.

Ninguna proporción de partición resultó mejor de forma consistente. Las particiones 90-10 evalúan sobre 18, 57 y 137 patrones, de modo que un solo error mueve la exactitud entre 0.7 y 5.6 puntos, como el cien por ciento del Perceptrón en Wine 90-10, que equivale a un vino más acertado que el MLP.

% ============================================================
\section{Referencias}
\begin{itemize}
    \item Rumelhart, D. E., Hinton, G. E., \& Williams, R. J. (1986). Learning representations by back-propagating errors. \textit{Nature}, 323(6088), 533--536.
    \item Haykin, S. (2009). \textit{Neural Networks and Learning Machines} (3.\textsuperscript{a} ed.). Pearson.
    \item Redes neuronales, algoritmo Backpropagation (2007). Guía de laboratorio, material del curso Inteligencia Computacional Aplicada, Universidad Distrital Francisco José de Caldas.
    \item Fisher, R. A. (1936). The use of multiple measurements in taxonomic problems. \textit{Annals of Eugenics}, 7(2), 179--188.
    \item Aeberhard, S., \& Forina, M. (1991). Wine. UCI Machine Learning Repository.
    \item Wolberg, W., Mangasarian, O., Street, N., \& Street, W. (1993). Breast Cancer Wisconsin (Diagnostic). UCI Machine Learning Repository.
    \item Lohweg, V. (2013). Banknote Authentication. UCI Machine Learning Repository.
\end{itemize}

% ============================================================
\appendix
\section{Tablas complementarias}
\label{anexo:tablas}

\begin{table}[!htb]
    \centering
    \small
    \caption{Matrices de confusión de las doce configuraciones sobre Iris. Cada celda es el número de flores de la clase real en fila asignadas a la clase en columna, con las clases en el orden setosa, versicolor, virginica}
    \label{tab:iris_todas}
    \begin{tabular}{c S[table-format=1.2] *{9}{S[table-format=2.0]}}
        \toprule
        \textbf{Ocultas} & {$\boldsymbol{\eta}$} & {\textbf{se-se}} & {\textbf{se-ve}} & {\textbf{se-vi}} & {\textbf{ve-se}} & {\textbf{ve-ve}} & {\textbf{ve-vi}} & {\textbf{vi-se}} & {\textbf{vi-ve}} & {\textbf{vi-vi}} \\
        \midrule
        3     & 0.05 & 14 & 0 & 0 & 0 & 11 & 2 & 0 & 0 & 18 \\
        3     & 0.1  & 14 & 0 & 0 & 0 & 12 & 1 & 0 & 0 & 18 \\
        3     & 0.3  & 14 & 0 & 0 & 0 & 12 & 1 & 0 & 0 & 18 \\
        6     & 0.05 & 14 & 0 & 0 & 0 & 11 & 2 & 0 & 0 & 18 \\
        6     & 0.1  & 14 & 0 & 0 & 0 & 12 & 1 & 0 & 0 & 18 \\
        6     & 0.3  & 14 & 0 & 0 & 0 & 12 & 1 & 0 & 0 & 18 \\
        12    & 0.05 & 14 & 0 & 0 & 0 & 11 & 2 & 0 & 0 & 18 \\
        12    & 0.1  & 14 & 0 & 0 & 0 & 12 & 1 & 0 & 0 & 18 \\
        12    & 0.3  & 14 & 0 & 0 & 0 & 12 & 1 & 0 & 0 & 18 \\
        {[8 4]} & 0.05 & 14 & 0 & 0 & 0 & 12 & 1 & 0 & 0 & 18 \\
        {[8 4]} & 0.1  & 14 & 0 & 0 & 0 & 12 & 1 & 0 & 0 & 18 \\
        {[8 4]} & 0.3  & 14 & 0 & 0 & 0 & 12 & 1 & 0 & 0 & 18 \\
        \bottomrule
    \end{tabular}
\end{table}

\lstset{basicstyle=\ttfamily\scriptsize}
\section{Código de los bloques completados}
\label{anexo:codigo}

Se reproducen completos los cuatro bloques que la plantilla dejaba pendientes y la función de entrenamiento con parada temprana que se agregó para los experimentos. Ningún nombre ni argumento de las funciones de la plantilla fue modificado.

\noindent\begin{minipage}{\linewidth}
\lstinputlisting[language=Matlab,caption={Inicialización de pesos y sesgos}]{fragmentos/init_mlp.m}
\end{minipage}
\vspace{0.6em}

\noindent\begin{minipage}{\linewidth}
\lstinputlisting[language=Matlab,caption={Propagación hacia adelante}]{fragmentos/forward_mlp.m}
\end{minipage}
\vspace{0.6em}

\noindent\begin{minipage}{\linewidth}
\lstinputlisting[language=Matlab,caption={Retropropagación del error y actualización con momento}]{fragmentos/backward_mlp.m}
\end{minipage}
\vspace{0.6em}

\noindent\begin{minipage}{\linewidth}
\lstinputlisting[language=Matlab,caption={Derivadas de las funciones de activación}]{fragmentos/activation_deriv.m}
\end{minipage}
\vspace{0.6em}

\noindent\begin{minipage}{\linewidth}
\lstinputlisting[language=Matlab,caption={Entrenamiento con error de validación y parada por épocas sin mejora}]{fragmentos/train_mlp.m}
\end{minipage}

\end{document}
